% =====================================================================
%  6. Protocole d'évaluation
% =====================================================================
\clearpage
\section{Protocole de validation et métriques}
\label{sec:protocole}

\subsection{Pourquoi une validation croisée imbriquée ?}

Un modèle doit être jugé sur des données qu'il n'a pas vues ; c'est le rôle de la
validation croisée~\cite{sklearn_cv}. Mais lorsqu'on \emph{choisit} aussi des
hyperparamètres, voire une famille de modèles, en regardant des scores de
validation, ces scores deviennent optimistes : on a sélectionné le meilleur
résultat parmi beaucoup d'essais~\cite{varma_simon_2006,cawley_talbot_2010}.

La validation croisée \textbf{imbriquée} sépare les deux rôles :
\begin{itemize}
  \item une \textbf{boucle interne} (3 plis) choisit les hyperparamètres --- et, pour la \og sélection interne \fg{}, la famille --- en n'utilisant que le train externe ;
  \item une \textbf{boucle externe} (5 plis) évalue \emph{une seule fois} la configuration choisie sur des données jamais utilisées pour la choisir.
\end{itemize}
Le score externe estime donc la performance de \emph{toute la procédure} de
réglage, et non celle d'une configuration choisie après coup.

\begin{figure}[htbp]
\centering
\begin{tikzpicture}[x=1cm,y=1cm,
  tr/.style={draw=white, line width=0.8pt, fill=bleuviz!75},
  va/.style={draw=white, line width=0.8pt, fill=orangeviz!90},
  itr/.style={draw=white, line width=0.6pt, fill=bleuviz!40},
  iva/.style={draw=white, line width=0.6pt, fill=vertviz!80},
  lab/.style={font=\scriptsize, color=encre2, anchor=east}]
  % Boucle externe
  \foreach \k in {1,...,5} {
    \pgfmathsetmacro{\y}{-0.5*(\k-1)}
    \node[lab] at (-0.1,\y+0.2) {Pli externe \k};
    \foreach \j in {1,...,5} {
      \pgfmathsetmacro{\x}{1.6*(\j-1)}
      \ifnum\j=\k
        \fill[va] (\x,\y) rectangle ++(1.6,0.4);
      \else
        \fill[tr] (\x,\y) rectangle ++(1.6,0.4);
      \fi
    }
  }
  % Zoom sur le train externe du pli 1
  \draw[decorate, decoration={brace, amplitude=5pt, mirror}, color=csgris, line width=0.7pt]
    (1.6,-2.15) -- (8.0,-2.15) node[midway, below=6pt, font=\scriptsize, color=encre2]
    {train externe du pli 1, redécoupé en 3 plis internes};
  \foreach \k in {1,...,3} {
    \pgfmathsetmacro{\y}{-3.2-0.42*(\k-1)}
    \node[lab] at (1.5,\y+0.17) {interne \k};
    \foreach \j in {1,...,3} {
      \pgfmathsetmacro{\x}{1.6+2.1333*(\j-1)}
      \ifnum\j=\k
        \fill[iva] (\x,\y) rectangle ++(2.1333,0.34);
      \else
        \fill[itr] (\x,\y) rectangle ++(2.1333,0.34);
      \fi
    }
  }
  % Légende et étapes
  \node[font=\scriptsize, align=left, anchor=north west, text width=4.3cm] at (8.4,0.4) {
    \tikz\fill[bleuviz!75] (0,0) rectangle (0.3,0.2); \ train externe\\[2pt]
    \tikz\fill[orangeviz!90] (0,0) rectangle (0.3,0.2); \ validation externe (évaluation)\\[2pt]
    \tikz\fill[bleuviz!40] (0,0) rectangle (0.3,0.2); \ train interne\\[2pt]
    \tikz\fill[vertviz!80] (0,0) rectangle (0.3,0.2); \ validation interne (réglage)};
  \node[font=\scriptsize, align=left, anchor=north west, text width=4.3cm] at (8.4,-2.0) {
    \textbf{Pour chaque pli externe :}\\
    1. 39 configurations $\times$ 3 plis internes ;\\
    2. choix par MAE interne moyenne ;\\
    3. réapprentissage sur tout le train externe ;\\
    4. prédiction unique de la validation externe.};
\end{tikzpicture}
\caption{Validation croisée imbriquée : 5 plis externes pour l'évaluation, 3 plis
internes dans chaque train externe pour le réglage. Le prétraitement est réappris
à chaque ajustement (interne et externe).}
\label{fig:cv_imbriquee}
\end{figure}

Pour chaque cible et chaque pli externe, cela représente 39 configurations
$\times$ 3 validations internes, plus un réapprentissage par famille : dans le
benchmark du notebook 03, \textbf{2\,340 ajustements internes et 100 réapprentissages
externes}, soit 2\,440 ajustements d'estimateurs. Tout échec d'ajustement lève une
erreur (\texttt{error\_score='raise'}) plutôt que d'être masqué.

\subsection{Deux scénarios de validation, deux questions}
\label{subsec:deux_scenarios}

La manière de former les plis détermine la question à laquelle répond le score.
Nous utilisons deux scénarios, avec les mêmes labels et les mêmes familles
(\cref{tab:scenarios}).

\begin{table}[htbp]
\centering
\small
\caption{Les deux scénarios de validation croisée.}
\label{tab:scenarios}
\begin{tabularx}{\textwidth}{@{}>{\raggedright\arraybackslash}p{2.9cm} L L@{}}
\toprule
 & \textbf{Validation aléatoire} & \textbf{Validation groupée} \\
\midrule
Découpage externe & KFold(5) mélangé, graine 42, par cible & GroupKFold(5) sur les séries inférées, sans mélange \\
Découpage interne & KFold(3) mélangé, graine $43 + \text{pli}$ & GroupKFold(3) dans chaque train externe \\
Ce qui est séparé & Des lignes & Des séries entières \\
Question posée & Prédire une nouvelle soudure d'une série \emph{déjà représentée} & Prédire les soudures d'une \emph{nouvelle étude} \\
Partage observé & 99 à 100\,\% des lignes de validation ont une série présente dans le train & Aucune série partagée \\
Risque & Optimisme si l'usage visé concerne de nouvelles séries & Pessimisme si les séries sont très hétérogènes ou mal définies \\
\bottomrule
\end{tabularx}
\end{table}

\begin{intuition}
Imaginons une étude qui fait varier la teneur en nickel de cinq soudures
fabriquées avec le même procédé, dans le même laboratoire. En validation
aléatoire, quatre d'entre elles peuvent être dans le train : le modèle
\og connaît \fg{} déjà le niveau de cette étude et n'a qu'à interpoler. En
validation groupée, les cinq sont absentes : le modèle doit deviner le niveau
d'une étude qu'il n'a jamais vue. Les données structurées en groupes appellent
une validation qui respecte cette structure~\cite{roberts_2017}.
\end{intuition}

En validation groupée, chaque pli externe contient 8 à 10 groupes ; les
validations des trois cibles de traction comportent 103 à 166 labels selon le
pli. Le nombre de lignes ne garantit pas le nombre d'observations réellement
indépendantes.

\begin{decision}
Les deux scénarios sont conservés et rapportés côte à côte. Le choix du
protocole \og principal \fg{} dépend de l'usage : la validation aléatoire décrit
l'interpolation dans la population couverte par la base, la validation groupée
est un test de transfert à de nouvelles études. On ne choisit pas le protocole
qui affiche le meilleur score.
\end{decision}

\subsection{Métriques}
\label{subsec:metriques}

Pour $n$ labels de validation $y_i$, de prédictions $\hat{y}_i$ et de moyenne $\bar{y}$ :
\begin{align}
  \mathrm{MAE}  &= \frac{1}{n}\sum_{i=1}^{n} |y_i - \hat{y}_i|, \\
  \mathrm{RMSE} &= \sqrt{\frac{1}{n}\sum_{i=1}^{n} (y_i - \hat{y}_i)^2}, \\
  R^2 &= 1 - \frac{\sum_{i=1}^{n} (y_i - \hat{y}_i)^2}{\sum_{i=1}^{n} (y_i - \bar{y})^2}.
\end{align}

\begin{itemize}
  \item La \textbf{MAE} est l'erreur typique, dans l'unité de la cible (MPa, points de \%, J). Elle sert de \textbf{critère de réglage}.
  \item La \textbf{RMSE} pénalise davantage les grandes erreurs ; l'écart entre RMSE et MAE renseigne sur la présence d'erreurs très fortes.
  \item Le \textbf{$R^2$} est sans unité : c'est la part de la variance de la validation expliquée par le modèle. $R^2 = 0$ correspond à prédire la moyenne \emph{de la validation} ; $R^2 < 0$ signifie faire pire. Il permet de comparer des cibles d'unités différentes, mais dépend de la variance de la population évaluée.
\end{itemize}

Deux agrégations coexistent. La \textbf{moyenne des plis} calcule la métrique dans
chaque pli externe puis en fait la moyenne non pondérée (avec l'écart-type SD) ;
le score \textbf{OOF} (\emph{out-of-fold}) réunit toutes les prédictions externes
puis calcule une seule métrique. Les deux diffèrent, surtout pour le $R^2$, car
chaque pli a sa propre variance. Sauf mention contraire, les tableaux donnent la
moyenne des plis. Pour comparer deux méthodes sur les mêmes observations, nous
utilisons un \textbf{gain apparié} de MAE OOF accompagné d'un intervalle
\textbf{bootstrap par groupes} (2\,000 rééchantillonnages des séries)~\cite{efron_tibshirani_1993},
plus honnête qu'un bootstrap par lignes lorsque les lignes d'une série sont dépendantes.

\begin{precaution}
Les écarts-types entre plis et les intervalles bootstrap sont descriptifs : ils
sont conditionnels aux plis et aux modèles entraînés, et ne constituent pas des
intervalles de confiance sur de nouvelles publications. Un $R^2$ de 0,5 n'est
pas un seuil garanti par le choix d'un algorithme.
\end{precaution}

\subsection{Protocole inductif pour le semi-supervisé}
\label{subsec:inductif}

Un protocole semi-supervisé est \emph{transductif} s'il exploite les entrées de
validation (sans leurs labels) pendant l'apprentissage, et \emph{inductif} s'il
ne les voit jamais. Nous adoptons un protocole \textbf{strictement inductif} :
\begin{itemize}
  \item le pool non étiqueté d'un pli ne contient que des lignes sans la cible \emph{affectées au train courant} --- interne comme externe ;
  \item en validation aléatoire, les lignes sans label sont réparties en cinq blocs équilibrés (graine 142), puis en trois blocs internes (graine $1042 + \text{pli}$) ; aucune ligne de validation n'entre dans le pool ;
  \item en validation groupée, le pool exclut en plus \emph{toutes les séries} de validation ;
  \item ni les entrées ni les labels de validation ne servent au calcul de confiance, à l'arrêt ou au prétraitement ;
  \item les modèles réentraînés sur tous les labels ne servent jamais à scorer les plis.
\end{itemize}

\subsection{Graines et reproductibilité}

Toutes les sources d'aléa sont fixées et enregistrées dans les manifestes
d'expérience (\cref{tab:graines}). Les partitions sont exportées par identifiant
brut ; les versions des bibliothèques et les empreintes SHA-256 des données, du
code et des exports sont consignées (\cref{annexe:reproductibilite}).

\begin{table}[htbp]
\centering
\small
\caption{Graines aléatoires utilisées.}
\label{tab:graines}
\begin{tabularx}{\textwidth}{@{}l L@{}}
\toprule
Graine & Usage \\
\midrule
42 & Estimateurs stochastiques (forêts), imputation itérative, plis externes aléatoires, base des graines de COREG, bootstrap des gains COREG (2\,000 tirages), test de permutation série/cible (999) \\
$43 + \text{pli}$ & Plis internes de la validation aléatoire \\
142 & Répartition des lignes sans label en cinq blocs (validation aléatoire) \\
$1042 + \text{pli}$ & Blocs internes des lignes sans label \\
1042 & CV5 aléatoire de sélection des modèles réentraînés sur tous les labels \\
1142 & Affectation des lignes sans label lors des réapprentissages complets de COREG \\
43, 44 & Sensibilité de COREG à la graine (à $k$ fixé) \\
\bottomrule
\end{tabularx}
\end{table}

\begin{aretenir}
Validation croisée \textbf{imbriquée} 5$\times$3, prétraitement réappris dans
chaque train, sélection par MAE interne. Deux scénarios complémentaires :
\textbf{aléatoire} (interpolation dans des séries connues) et \textbf{groupé}
(transfert à de nouvelles séries). Métriques MAE, RMSE, $R^2$ ; semi-supervisé
évalué selon un protocole \textbf{inductif}.
\end{aretenir}
